\documentclass[12pt]{article}
\usepackage[utf8]{inputenc}
\usepackage[T1]{fontenc}
\usepackage{hyperref}
\title{Copainal\'a Zoque Language Resources for K-12 Curriculum Development}
\author{Compiled by Claude Sonnet 5.5 (claude-sonnet-5-5) for Andr\'as Kornai}
\date{}
\begin{document}
\maketitle

\section{Introduction}
Copainal\'a Zoque (ISO 639-3 zoc) is one of the Chiapas Zoque varieties of the Mixe-Zoquean family, spoken in and around Copainal\'a, Chiapas, Mexico, by roughly 10,000 people. The main resources found are older SIL linguistic works. No school curriculum materials were found.

\section{Primary Educational Resources}
No primers, textbooks or school programmes specific to Copainal\'a Zoque were found in search results.

\section{Linguistic and Reference Materials}
\begin{itemize}
\item Harrison, Harrison and Garc\'ia Hern\'andez, \emph{Diccionario zoque de Copainal\'a} (Instituto Ling\"u\'istico de Verano, 1981), library catalogue record: \url{https://libcat.library.tamu.edu/Record/in00000725180}
\item Harrison and Harrison 1984, \emph{Survey of morphology and syntax for Zoque of Copainala} (SIL Mexico Workpapers 6), Glottolog reference: \url{https://glottolog.org/resource/reference/id/76152}
\item Harrison and Harrison 1948, Spanish--Zoque dictionary, Glottolog reference: \url{https://glottolog.org/resource/reference/id/81207}
\item Colegio de M\'exico repository record (returned for this variety; contents not confirmed): \url{https://repositorio.colmex.mx/concern/books/xs55mg75s?locale=en}
\item Wikipedia, Chiapas Zoque: \url{https://en.wikipedia.org/wiki/Chiapas_Zoque}
\item Zoque materials in the C.~F. Voegelin Papers (American Philosophical Society): \url{https://indigenousguide.amphilsoc.org/entry/40326}
\item Crossgram language page for Copainal\'a Zoque: \url{https://crossgram.clld.org/languages/copa1236}
\end{itemize}

\section{Religious and Community Resources}
None specific to this variety were found.

\section{Digital and Audio Resources}
No audio or video resources specific to Copainal\'a Zoque were found. The sources found are catalogue records and bibliographic entries, and no full digitized text was seen.

\section{Curriculum Development Context}
The language has a sizeable community and a published dictionary and grammar sketches, which gives a starting point for a curriculum. Those works are decades old, however, and mostly lack open digital access. A curriculum would need new pedagogical materials produced with Copainal\'a speakers, and access to the SIL dictionary through libraries or archives.

\end{document}